\frfilename{mt276.tex}
\versiondate{16.4.13}
\copyrightdate{2000}

\def\chaptername{Teori probabilitas}
\def\sectionname{Barisan selisih martingal}

\newsection{276}

Berdampingan dengan konsep `martingal' terdapat konsep `barisan selisih
martingal' (276A), suatu perumuman langsung dari gagasan `barisan bebas'.
Dalam bagian ini saya mengumpulkan hasil-hasil yang dapat dinyatakan secara
alami dengan barisan selisih, termasuk satu lagi hukum kuat bilangan besar
(276C). Saya mengakhiri bagian ini dengan bukti teorema Koml\'os (276H).

\leader{276A}{Barisan selisih martingal (a)}\cmmnt{ Jika
$\sequencen{X_n}$
adalah martingal yang disesuaikan dengan suatu barisan
$\sequencen{\Sigma_n}$ aljabar-sigma, maka

\Centerline{$\int_EX_{n+1}-X_n=0$}

\noindent kapan pun $E\in\Sigma_n$.} Kita katakan bahwa jika
$(\Omega,\Sigma,\mu)$ adalah ruang probabilitas, dengan pelengkapan
$(\Omega,\hat\Sigma,\hat\mu)$, dan $\sequencen{\Sigma_n}$ adalah
barisan tak menurun subaljabar-sigma dari $\hat\Sigma$, maka suatu
{\bf barisan selisih martingal yang disesuaikan dengan
$\sequencen{\Sigma_n}$} adalah barisan $\sequencen{X_n}$ variabel acak
bernilai real pada $\Omega$, semuanya berekspektasi hingga, sedemikian
sehingga (i) $\dom X_n\in\Sigma_n$ dan $X_n$ terukur terhadap
$\Sigma_n$, untuk setiap $n\in\Bbb N$ (ii) $\int_EX_{n+1}=0$ kapan pun
$n\in\Bbb N$ dan $E\in\Sigma_n$.

\header{276Ab}{\bf (b)}
\cmmnt{ Jelas bahwa} $\sequencen{X_n}$ adalah barisan selisih martingal
yang disesuaikan dengan $\sequencen{\Sigma_n}$ jika dan hanya jika
$\sequencen{\sum_{i=0}^nX_i}$ adalah martingal yang disesuaikan dengan
$\sequencen{\Sigma_n}$.

\header{276Ac}{\bf (c)}
Seperti pada 275Cd, kita dapat mengatakan bahwa suatu barisan
$\sequencen{X_n}$ dengan sendirinya merupakan {\bf barisan selisih
martingal} jika $\sequencen{\sum_{i=0}^nX_i}$ adalah martingal\cmmnt{,
yaitu jika $\sequencen{X_n}$ adalah barisan selisih martingal yang
disesuaikan dengan $\sequencen{\tilde\Sigma_n}$, dengan
$\tilde\Sigma_n$ aljabar-sigma yang dibangkitkan oleh
$\bigcup_{i\le n}\Sigma_{X_i}$}.

\header{276Ad}{\bf (d)} Jika $\sequencen{X_n}$ adalah barisan selisih
martingal, maka $\sequencen{a_nX_n}$ adalah barisan selisih martingal
untuk sebarang bilangan real $a_n$.

\header{276Ae}{\bf (e)} Jika $\sequencen{X_n}$ adalah barisan selisih
martingal dan $X'_n\eae X_n$ untuk setiap $n$, maka
$\sequencen{X'_n}$ adalah barisan selisih martingal.
\cmmnt{(Bandingkan 275Ce.)}

\cmmnt{\header{276Af}{\bf (f)} Tentu saja contoh terpenting dari
`barisan selisih martingal' adalah contoh dalam 275Bb: setiap barisan
bebas variabel acak dengan ekspektasi nol merupakan barisan selisih
martingal. Ternyata beberapa teorema dalam \S273 mengenai barisan bebas
semacam itu dapat diperumum ke barisan selisih martingal.}

\leader{276B}{Proposisi} Misalkan $\sequencen{X_n}$ adalah barisan
selisih martingal sedemikian sehingga
$\sum_{n=0}^{\infty}\Expn(X_n^2)<\infty$. Maka
$\sum_{n=0}^{\infty}X_n$ terdefinisi dan berhingga hampir di mana-mana.

\proof{{\bf (a)} Misalkan $(\Omega,\Sigma,\mu)$ adalah ruang
probabilitas yang mendasari, $(\Omega,\hat\Sigma,\hat\mu)$
pelengkapannya, dan $\sequencen{\Sigma_n}$ suatu barisan tak menurun
subaljabar-sigma dari $\hat\Sigma$ sedemikian sehingga
$\sequencen{X_n}$ disesuaikan dengan $\sequencen{\Sigma_n}$. Tetapkan
$Y_n=\sum_{i=0}^nX_i$ untuk setiap $n\in\Bbb N$. Maka
$\sequencen{Y_n}$ adalah martingal yang disesuaikan dengan
$\sequencen{\Sigma_n}$.

\medskip

{\bf (b)} $\Expn(Y_{n}\times X_{n+1})=0$ untuk setiap $n$. \Prf\ $Y_n$
adalah jumlah variabel acak dengan varians hingga, sehingga
$\Expn(Y_n^2)<\infty$, menurut 244Ba; akibatnya
$Y_n\times X_{n+1}$ berekspektasi hingga, menurut 244Eb. Karena fungsi
konstan $\pmb{0}$ merupakan ekspektasi bersyarat dari $X_{n+1}$ pada
$\Sigma_n$,

\Centerline{$\Expn(Y_n\times X_{n+1})=\Expn(Y_n\times\pmb{0})=0$,}

\noindent menurut 242L.\ \Qed

\medskip

{\bf (c)} Akibatnya $\Expn(Y_n^2)=\sum_{i=0}^n\Expn(X_i^2)$ untuk
setiap $n$. \Prf\ Lakukan induksi pada $n$. Untuk langkah induksi,
kita mempunyai

\Centerline{$\Expn(Y_{n+1}^2)=\Expn(Y_n^2+2Y_n\times X_{n+1}+X_{n+1}^2)
=\Expn(Y_n^2)+\Expn(X_{n+1}^2)$}

\noindent karena, menurut (b), $\Expn(Y_n\times X_{n+1})=0$.\ \Qed

\medskip

{\bf (d)} Tentu saja

\Centerline{$\Expn(|Y_n|)
=\int|Y_n|\times\chi\Omega\le\|Y_n\|_2\|\chi\Omega\|_2
=\sqrt{\Expn(Y_n^2)}$,}

\noindent sehingga

\Centerline{$\sup_{n\in\Bbb N}\Expn(|Y_n|)
\le\sup_{n\in\Bbb N}\sqrt{\Expn(Y_n^2)}
=\sqrt{\sumop_{i=0}^{\infty}\Expn(X_i^2)}<\infty$.}

\noindent Menurut 275G, $\lim_{n\to\infty}Y_n$ terdefinisi dan
berhingga hampir di mana-mana, yaitu $\sum_{i=0}^{\infty}X_i$
terdefinisi dan berhingga hampir di mana-mana.
}%end of proof of 276B

\leader{276C}{Hukum kuat bilangan besar: bentuk keempat} Misalkan
$\sequencen{X_n}$ adalah barisan selisih martingal, dan andaikan bahwa
$\sequencen{b_n}$ adalah barisan tak menurun dalam $\ooint{0,\infty}$
yang divergen ke $\infty$, sedemikian sehingga
$\sum_{n=0}^{\infty}\Bover1{b_n^2}\Var(X_n)<\infty$. Maka

\Centerline{$\lim_{n\to\infty}\Bover{1}{b_n}
\sum_{i=0}^n X_i=0$}

\noindent hampir di mana-mana.

\proof{ (Bandingkan 273D.) Seperti biasa, tuliskan
$(\Omega,\Sigma,\mu)$ untuk ruang probabilitas yang mendasari. Tetapkan

\Centerline{$\tilde X_n=\Bover1{b_n}X_n$}

\noindent untuk setiap $n$; maka $\sequencen{\tilde X_n}$ juga merupakan
barisan selisih martingal, dan

\Centerline{$\sum_{n=1}^{\infty}\Expn(\tilde X_n^2)
=\sum_{n=1}^{\infty}\Bover1{b_n^2}\Var(X_n)<\infty$.}

\noindent Menurut 276B, $\sequencen{\tilde X_n(\omega)}$ dapat dijumlahkan
untuk hampir setiap $\omega\in\Omega$. Tetapi menurut 273Cb,

\Centerline{$\lim_{n\to\infty}\Bover{1}{b_n}\sum_{i=0}^n X_i(\omega)
=\lim_{n\to\infty}\Bover{1}{b_n}\sum_{i=0}^nb_i\tilde X_i(\omega)
=0$}

\noindent untuk setiap $\omega$ semacam itu. Jadi hasilnya terbukti.
}%end of proof of 276C

\leader{276D}{Akibat} Misalkan $\sequencen{X_n}$ adalah martingal
sedemikian sehingga $b_n=\Expn(X_n^2)$ berhingga untuk setiap $n$.

(a) Jika $\sup_{n\in\Bbb N}b_n$ tak berhingga, maka
$\lim_{n\to\infty}\bover1{b_n}X_n=0$ hampir di mana-mana.

(b) Jika $\sup_{n\ge 1}\bover1nb_n<\infty$, maka
$\lim_{n\to\infty}\bover1nX_n=0$ hampir di mana-mana.

\proof{ Pertimbangkan barisan selisih martingal
$\sequencen{Y_n}=\sequencen{X_{n+1}-X_n}$. Maka
$\Expn(Y_n\times X_n)=0$, sehingga
$\Expn(Y_n^2)+\Expn(X_{n}^2)=\Expn(X_{n+1}^2)$ untuk setiap $n$.
Khususnya, $\sequencen{b_n}$ harus tak menurun.

\medskip

{\bf (a)} Jika $\lim_{n\to\infty}b_n=\infty$, ambil $m$ sedemikian
sehingga $b_m>0$; maka

\Centerline{$\sum_{n=m}^{\infty}\Bover1{b_{n+1}^2}\Var(Y_n)
=\sum_{n=m}^{\infty}\Bover1{b_{n+1}^2}(b_{n+1}-b_n)
\le\int_{b_m}^{\infty}\Bover1{t^2}dt<\infty$.}

\noindent Menurut 276C (dengan mengubah $b_i$ untuk $i<m$ jika perlu),

\Centerline{$\lim_{n\to\infty}\Bover1{b_n}X_n
=\lim_{n\to\infty}\Bover1{b_{n+1}}(X_0+\sum_{i=0}^nY_i)
=\lim_{n\to\infty}\Bover1{b_{n+1}}\sum_{i=0}^nY_i=0$}

\noindent hampir di mana-mana.

\medskip

{\bf (b)} Jika $\gamma=\sup_{n\ge 1}\bover1nb_n<\infty$, maka
$\bover1{(n+1)^2}\le\min(1,\gamma^2/t^2)$ untuk $b_n<t\le b_{n+1}$,
sehingga

\Centerline{$\sum_{n=0}^{\infty}\Bover1{(n+1)^2}(b_{n+1}-b_n)
\le \gamma+\gamma^2\int_{\gamma}^{\infty}\Bover1{t^2}dt<\infty$,}

\noindent dan, dengan argumen yang sama seperti sebelumnya,
$\lim_{n\to\infty}\bover1nX_n=0$ hampir di mana-mana.

}%end of proof of 276D

\leader{276E}{`Ketidakmungkinan sistem'}\cmmnt{ {\bf (a)} Saya kembali
ke kata `martingal' dan gagasan sistem perjudian. Pertimbangkan seorang
penjudi yang mengambil serangkaian taruhan `adil', yaitu taruhan dengan
ekspektasi hasil nol, tetapi memilih taruhan mana yang akan diambil
berdasarkan pengalaman sebelumnya. Model yang sesuai untuk barisan
kejadian acak semacam itu adalah martingal dalam pengertian 275A, dengan
$\Sigma_n$ sebagai aljabar semua kejadian yang dapat diamati sampai dengan
hasil taruhan ke-$n$, dan $X_n$ sebagai keuntungan bersih penjudi pada saat
itu. (Dalam model ini wajar untuk mengambil
$\Sigma_0=\{\emptyset,\Omega\}$ dan $X_0=\tbf{0}$.) Paradoks tertentu dapat
muncul jika kita mencoba membayangkan model ini dengan $\Sigma_n$ tanpa
atom; sebagai permulaan mungkin lebih mudah bekerja dengan kasus diskret,
di mana setiap $\Sigma_n$ berhingga, atau merupakan himpunan gabungan dari
suatu keluarga kejadian atomik terhitung. Sekarang andaikan taruhan yang
bersangkutan hanyalah taruhan dua-arah, dengan dua hasil yang sama
mungkin, tetapi penjudi memilih besar taruhannya setiap kali. Dalam kasus
ini kita dapat memandang hasil-hasil tersebut sebagai bersesuaian dengan
barisan bebas $\sequencen{W_n}$ variabel acak, yang masing-masing mengambil
nilai $\pm 1$ dengan probabilitas sama. Sistem penjudi itu harus berbentuk

\Centerline{$X_{n+1}=X_n + Z_{n+1}\times W_{n+1}$,}

\noindent dengan $Z_{n+1}$ besar taruhannya pada taruhan ke-$(n+1)$, dan
harus konstan pada setiap atom aljabar-sigma $\Sigma_n$ yang dibangkitkan
oleh $W_1,\ldots,W_n$. Intinya, karena $\int_EW_{n+1}=0$ untuk setiap
$E\in\Sigma_n$, $\Expn(Z_{n+1}\times W_{n+1})=0$, sehingga
$\Expn(X_{n+1})=\Expn(X_n)$.

\header{276Eb}} %end of comment
{\bf (b)}\cmmnt{ Hasil umum yang memuat ini sebagai kasus khusus adalah
sebagai berikut.} Jika $\sequencen{W_n}$ adalah barisan selisih martingal
yang disesuaikan dengan $\sequencen{\Sigma_n}$, dan
$\langle{Z_n}\rangle_{n\ge 1}$ adalah barisan variabel acak sedemikian
sehingga (i) $Z_n$ terukur terhadap $\Sigma_{n-1}$ (ii)
$Z_n\times W_{n}$ berekspektasi hingga untuk setiap $n\ge 1$, maka
$W_0,Z_1\times W_1,Z_2\times W_2,\ldots$ adalah barisan selisih martingal
yang disesuaikan dengan $\sequencen{\Sigma_n}$\cmmnt{; bukti bahwa
$\int_EZ_{n+1}\times W_{n+1}=0$ untuk setiap $E\in\Sigma_n$ persis
merupakan argumen dalam bagian (b) bukti 276B}.

\cmmnt{
\header{276Ec}{\bf (c)} Saya tadi mengajak Anda membatasi gagasan Anda
sejenak pada kasus diskret; tetapi jika Anda merasa memahami apa artinya
mengatakan bahwa suatu `sistem' atau {\bf barisan terprediksi}
$\langle{Z_n}\rangle_{n\ge 1}$ harus disesuaikan dengan
$\sequencen{\Sigma_n}$, dalam arti bahwa setiap $Z_n$ terukur terhadap
$\Sigma_{n-1}$, maka kesulitan selanjutnya terletak pada teori ukuran yang
diperlukan untuk menunjukkan bahwa integral
$\int_EZ_{n+1}\times W_{n+1}$ bernilai nol, dan itulah pokok buku ini.

\header{276Ed}{\bf (d)} Pertimbangkan sistem perjudian yang disebutkan
dalam 275Cf. Di sini gagasannya adalah $W_n=\pm 1$, seperti pada (a), dan
$Z_{n+1}=2^na$ jika $X_n\le 0$, serta $0$ jika $X_n>0$; yaitu, penjudi
menggandakan besar taruhannya setiap kali sampai ia menang, lalu berhenti.
Tentu saja ia hampir pasti akhirnya menang, sehingga
$\lim_{n\to\infty}X_n=a$ hampir di mana-mana, meskipun
$\Expn(X_n)=0$ untuk setiap $n$. Kita dapat menghitung distribusi $X_n$:
untuk $n\ge 1$ kita mempunyai $\Pr(X_n=a)=1-2^{-n}$,
$\Pr(X_n=-(2^n-1)a)=2^{-n}$. Jadi $\Expn(|X_n|)=(2-2^{-n+1})a$ dan
konvergensi hampir di mana-mana dari $X_n$ merupakan contoh Teorema
Konvergensi Martingal Doob.

Dalam bahasa waktu henti (275N), $X_n=\tilde Y_{\tau\wedge n}$,
dengan $Y_n=\sum_{k=0}^n2^kaW_k$ dan $\tau=\min\{n:Y_n>0\}$.
}

\leader{*276F}{}\cmmnt{ Sekarang saya beralih ke teorema Koml\'os.
Langkah pertama adalah penyempurnaan kecil dari 276C.

\medskip

\noindent}{\bf Lemma} Misalkan $(\Omega,\Sigma,\mu)$ adalah ruang
probabilitas, dan $\sequencen{\Sigma_n}$ suatu barisan tak menurun
subaljabar-sigma dari $\Sigma$. Andaikan bahwa $\sequencen{X_n}$ adalah
barisan variabel acak pada $\Omega$ sedemikian sehingga (i) $X_n$
terukur terhadap $\Sigma_n$ untuk setiap $n$ (ii)
$\sum_{n=0}^{\infty}\Bover1{(n+1)^2}\Expn(X_n^2)$ berhingga (iii)
$\lim_{n\to\infty}X'_n=0$ hampir di mana-mana, dengan $X_n'$ suatu
ekspektasi bersyarat dari $X_n$ pada $\Sigma_{n-1}$ untuk setiap $n\ge 1$.
Maka $\lim_{n\to\infty}\Bover1{n+1}\sum_{k=0}^nX_k=0$ hampir di mana-mana.

\proof{ Dengan melakukan penyesuaian yang sesuai pada suatu himpunan
terabaikan jika perlu, kita dapat mengandaikan bahwa $X'_n$ terukur
terhadap $\Sigma_{n-1}$ untuk $n\ge 1$ dan bahwa setiap $X_n$ dan $X'_n$
terdefinisi pada seluruh $\Omega$. Tetapkan $X'_0=X_0$ dan
$Y_n=X_n-X'_n$ untuk $n\in\Bbb N$. Maka $\sequencen{Y_n}$ adalah barisan
selisih martingal yang disesuaikan dengan $\sequencen{\Sigma_n}$. Selain
itu, $\Expn(Y_n^2)\le\Expn(X_n^2)$ untuk setiap $n$. \Prf\ Jika $n\ge 1$,
$X'_n$ terintegralkan-kuadrat (244M), dan
$\Expn(Y_n\times X'_n)=0$, seperti dalam bagian (b) bukti 276B. Sekarang

\Centerline{$\Expn(X_n^2)
=\Expn((Y_n+X'_n)^2)
=\Expn(Y_n^2)+2\Expn(Y_n\times X'_n)+\Expn(X'_n)^2
\ge\Expn(Y_n^2)$.  \Qed}

Ini berarti bahwa
$\sum_{n=0}^{\infty}\Bover1{(n+1)^2}\Expn(Y_n^2)$ harus berhingga.
Menurut 276C,
$\lim_{n\to\infty}\Bover1{n+1}\sum_{i=0}^nY_i=0$ hampir di mana-mana.
Tetapi menurut 273Ca kita juga mempunyai
$\lim_{n\to\infty}\Bover1{n+1}\sum_{i=0}^nX'_i=0$ kapan pun
$\lim_{n\to\infty}X'_n=0$, yang berlaku hampir di mana-mana. Jadi
$\lim_{n\to\infty}\Bover1{n+1}\sum_{i=0}^nX_i=0$ hampir di mana-mana.
}%end of proof of 276F

\leader{*276G}{Lemma} Misalkan $(\Omega,\Sigma,\mu)$ adalah ruang
probabilitas, dan $\sequencen{X_n}$ suatu barisan variabel acak pada
$\Omega$ sedemikian sehingga $\sup_{n\in\Bbb N}\Expn(|X_n|)$ berhingga.
Untuk $k\in\Bbb N$ dan $x\in\Bbb R$, tetapkan $F_k(x)=x$ jika
$|x|\le k$, dan $0$ jika tidak. Misalkan $\Cal F$ suatu ultrafilter pada
$\Bbb N$.

(a) Untuk setiap $k\in\Bbb N$ terdapat fungsi terukur
$Y_k:\Omega\to[-k,k]$ sedemikian sehingga
$\lim_{n\to\Cal F}\int_EF_k(X_n)=\int_EY_k$ untuk setiap $E\in\Sigma$.

(b) $\lim_{n\to\Cal F}\Expn((F_k(X_n)-Y_k)^2)
\le\lim_{n\to\Cal F}\Expn(F_k(X_n)^2)$ untuk setiap $k$.

(c) $Y=\lim_{k\to\infty}Y_k$ terdefinisi hampir di mana-mana dan
$\lim_{k\to\infty}\Expn(|Y-Y_k|)=0$.

\proof{{\bf (a)} Untuk setiap $k$, $|F_k(X_n)|\leae k\chi\Omega$ bagi
setiap $n$, sehingga $\{F_k(X_n):n\in\Bbb N\}$ terintegralkan secara
seragam, dan $\{F_k(X_n)^{\ssbullet}:n\in\Bbb N\}$ relatif kompak lemah
dalam $L^1=L^1(\mu)$ (247C). Karena itu
$v_k=\lim_{n\to\Cal F}F_k(X_n)^{\ssbullet}$ terdefinisi dalam $L^1$
untuk topologi lemah (2A3Se); ambil $Y_k:\Omega\to\Bbb R$ sebagai fungsi
terukur sedemikian sehingga $Y_k^{\ssbullet}=v_k$. Untuk sebarang
$E\in\Sigma$,

\Centerline{$\int_EY_k
=\int v_k\times(\chi E)^{\ssbullet}
=\lim_{n\to\Cal F}\int_EF_k(X_n)$.}

\noindent Khususnya,

\Centerline{$|\int_EY_k|\le\sup_{n\in\Bbb N}|\int_EF_k(X_n)|\le k\mu E$}

\noindent untuk setiap $E$, sehingga $\{\omega:Y_k(\omega)>k\}$ dan
$\{\omega:Y_k(\omega)<-k\}$ keduanya terabaikan; dengan mengubah $Y_k$
pada suatu himpunan terabaikan jika perlu, kita dapat mengandaikan bahwa
$|Y_k(\omega)|\le k$ untuk setiap $\omega\in\Omega$.

\medskip

{\bf (b)} Karena $Y_k$ terbatas,
$Y_k^{\ssbullet}\in L^{\infty}(\mu)$, dan

\Centerline{$\lim_{n\to\Cal F}\int F_k(X_n)\times Y_k
=\lim_{n\to\Cal F}\int F_k(X_n)^{\ssbullet}\times Y_k^{\ssbullet}
=\int Y_k^{\ssbullet}\times Y_k^{\ssbullet}
=\int Y_k^2$.}

\noindent Karena itu

$$\eqalignno{\lim_{n\to\Cal F}\int(F_k(X_n)-Y_k)^2
&=\lim_{n\to\Cal F}\int F_k(X_n)^2
  -2\lim_{n\to\Cal F}\int F_k(X_n)\times Y_k
  +\int Y_k^2\cr
&=\lim_{n\to\Cal F}\int F_k(X_n)^2-\int Y_k^2
\le\lim_{n\to\Cal F}\int F_k(X_n)^2.\cr}$$

\medskip

{\bf (c)} Tetapkan $W_0=Y_0=\tbf{0}$, dan $W_k=Y_k-Y_{k-1}$ untuk
$k\ge 1$. Maka
$\Expn(|W_k|)\le\lim_{n\to\Cal F}\Expn(|F_k(X_n)-F_{k-1}(X_n)|)$
untuk setiap $k\ge 1$. \Prf\ Tetapkan
$E=\{\omega:W_k(\omega)\ge 0\}$. Maka

$$\eqalign{\int_EW_k
&=\int_EY_k-\int_EY_{k-1}\cr
&=\lim_{n\to\Cal F}\int_EF_k(X_n)-\lim_{n\to\Cal F}\int_EF_{k-1}(X_n)\cr
&=\lim_{n\to\Cal F}\int_EF_k(X_n)-F_{k-1}(X_n)
\le\lim_{n\to\Cal F}\int_E|F_k(X_n)-F_{k-1}(X_n)|.\cr}$$

\noindent Demikian pula,

\Centerline{$|\int_{\Omega\setminus E}W_k|
\le\lim_{n\to\Cal F}\int_{\Omega\setminus E}|F_k(X_n)-F_{k-1}(X_n)|$.}

\noindent Jadi

\Centerline{$\Expn(|W_k|)
=\int_EW_k-\int_{\Omega\setminus E}W_k
\le\lim_{n\to\Cal F}\int|F_k(X_n)-F_{k-1}(X_n)|$.  \Qed}

Akibatnya $\sum_{k=0}^{\infty}\Expn(|W_k|)$ berhingga. \Prf\ Untuk
sebarang $m\ge 1$,

$$\eqalign{\sum_{k=0}^m\Expn(|W_k|)
&\le\sum_{k=1}^m\lim_{n\to\Cal F}\Expn(|F_k(X_n)-F_{k-1}(X_n)|)\cr
&=\lim_{n\to\Cal F}\Expn(\sum_{k=1}^m|F_k(X_n)-F_{k-1}(X_n)|)\cr
&=\lim_{n\to\Cal F}\Expn(|F_m(X_n)|)
\le\sup_{n\in\Bbb N}\Expn(|X_n|).\cr}$$

\noindent Jadi
$\sum_{k=0}^{\infty}\Expn(|W_k|)\le\sup_{n\in\Bbb N}\Expn(|X_n|)$
berhingga.\ \Qed

Menurut teorema B. Levi (123A),
$\lim_{m\to\infty}\sum_{k=0}^m|W_k|$ berhingga hampir di mana-mana,
sehingga

\Centerline{$Y=\lim_{m\to\infty}Y_m=\sum_{k=0}^{\infty}W_k$}

\noindent terdefinisi hampir di mana-mana; dan selain itu

\Centerline{$\Expn(|Y-Y_k|)
\le\lim_{m\to\infty}\Expn(\sum_{j=k+1}^m|W_j|)
\to 0$}

\noindent ketika $k\to\infty$.
}%end of proof of 276G

\leader{*276H}{Teorema Koml\'os}\cmmnt{ ({\smc Koml\'os 67})}
Misalkan $(\Omega,\Sigma,\mu)$ adalah sebarang ruang ukur, dan
$\sequencen{X_n}$ suatu barisan fungsi terintegralkan bernilai real pada
$\Omega$ sedemikian sehingga $\sup_{n\in\Bbb N}\int|X_n|$ berhingga.
Maka terdapat subbarisan $\sequencen{X'_n}$ dari $\sequencen{X_n}$ dan
fungsi terintegralkan $Y$ sedemikian sehingga
$Y\eae\lim_{n\to\infty}\Bover1{n+1}\sum_{i=0}^nX''_i$ kapan pun
$\sequencen{X''_n}$ merupakan subbarisan dari $\sequencen{X'_n}$.

\proof{ Karena baik hipotesis maupun kesimpulan tidak terpengaruh oleh
perubahan $X_n$ pada suatu himpunan terabaikan, kita dapat mengandaikan
di seluruh pembahasan bahwa setiap $X_n$ terukur dan terdefinisi pada
seluruh $\Omega$. Selain itu, sebagai permulaan (sampai akhir bagian (e)
di bawah), andaikan bahwa $\mu\Omega=1$. Seperti dalam 276G, tetapkan
$F_k(x)=x$ untuk $|x|\le k$, dan $0$ untuk $|x|>k$.

\medskip

{\bf (a)} Misalkan $\Cal F$ sebarang ultrafilter nonprinsipal pada
$\Bbb N$ (2A1O). Untuk $j\in\Bbb N$, tetapkan
$p_j=\lim_{n\to\Cal F}\Pr(|X_n|>j)$. Maka
$\sum_{j=0}^{\infty}p_j$ berhingga. \Prf\ Untuk sebarang $k\in\Bbb N$,

$$\eqalignno{\sum_{j=0}^kp_j
&=\sum_{j=0}^k\lim_{n\to\Cal F}\Pr(|X_n|>j)
=\lim_{n\to\Cal F}\sum_{j=0}^k\Pr(|X_n|>j)\cr
&\le\lim_{n\to\Cal F}(1+\int|X_n|)
\le 1+\sup_{n\in\Bbb N}\int|X_n|.\cr}$$

\noindent Jadi
$\sum_{j=0}^{\infty}p_j\le 1+\sup_{n\in\Bbb N}\int|X_n|$ berhingga.\ \Qed

Dengan menetapkan

\Centerline{$p'_j=p_j-p_{j+1}=\lim_{n\to\Cal F}\Pr(j<|X_n|\le j+1)$}

\noindent untuk setiap $j$, kita mempunyai

$$\eqalign{\sum_{j=0}^{\infty}(j+1)p'_j
&=\lim_{m\to\infty}
  \bigl(\sum_{j=0}^m(j+1)p_j-\sum_{j=0}^m(j+1)p_{j+1}\bigr)\cr
&=\lim_{m\to\infty}\sum_{j=0}^mp_j-(m+1)p_{m+1}
\le\sum_{j=0}^{\infty}p_j
<\infty.\cr}$$

Selanjutnya,

\Centerline{$\lim_{n\to\Cal F}\int F_k(X_n)^2
\le\sum_{j=0}^k(j+1)^2p'_j$}

\noindent untuk setiap $k$. \Prf\ Dengan menetapkan
$E_{jn}=\{\omega:j<|X_n(\omega)|\le j+1\}$ untuk
$j$, $n\in\Bbb N$, kita mempunyai
$F_k(X_n)^2\le\sum_{j=0}^k(j+1)^2\chi E_{jn}$, sehingga

\Centerline{$\lim_{n\to\Cal F}\int F_k(X_n)^2
\le\lim_{n\to\Cal F}\sum_{j=0}^k(j+1)^2\mu E_{jn}
=\sum_{j=0}^k(j+1)^2p'_j$.  \Qed}

\medskip

{\bf (b)} Definisikan $\sequence{k}{Y_k}$ dan
$Y\eae\lim_{k\to\infty}Y_k$ dari $\sequencen{X_n}$ dan $\Cal F$
seperti dalam Lemma 276G. Maka

\Centerline{$J_k
=\{n:n\in\Bbb N,\,\int(F_k(X_n)-Y_k)^2\le 1+\sum_{j=0}^k(j+1)^2p'_j\}$}

\noindent termasuk dalam $\Cal F$ untuk setiap $k\in\Bbb N$.
\Prf\ Menurut (a) di atas dan 276Gb,

\Centerline{$\lim_{n\to\Cal F}\int(F_k(X_n)-Y_k)^2
\le\lim_{n\to\Cal F}\int F_k(X_n)^2
\le\sum_{j=0}^k(j+1)^2p'_j$.  \Qed}

\noindent Selain itu, tentu saja,

\Centerline{$K_k=\{n:n\in\Bbb N,\,\Pr(F_j(X_n)\ne X_n)\le p_j+2^{-j}$
untuk setiap $j\le k\}$}

\noindent termasuk dalam $\Cal F$ untuk setiap $k$.

\medskip

{\bf (c)} Untuk $n$, $k\in\Bbb N$, misalkan $Z_{kn}$ suatu fungsi
sederhana sedemikian sehingga $|Z_{kn}|\le|F_k(X_n)-Y_k|$ dan
$\int|F_k(X_n)-Y_k-Z_{kn}|\le 2^{-k}$. Untuk $m\in\Bbb N$, misalkan
$\Sigma_m$ aljabar himpunan bagian dari $\Omega$ yang dibangkitkan oleh
himpunan berbentuk $\{\omega:Z_{kn}(\omega)=\alpha\}$ untuk $k$, $n\le m$
dan $\alpha\in\Bbb R$. Karena setiap $Z_{kn}$ hanya mengambil berhingga
banyak nilai, $\Sigma_m$ berhingga (dan oleh karena itu merupakan
subaljabar-sigma dari $\Sigma$); dan tentu saja
$\Sigma_m\subseteq\Sigma_{m+1}$ untuk setiap $m$.

Kita perlu menelaah ekspektasi bersyarat pada $\Sigma_m$, dan karena
$\Sigma_m$ selalu berhingga, ekspektasi ini memiliki bentuk yang sangat
langsung. Misalkan $\Cal A_m$ himpunan `atom', atau himpunan tak kosong
minimal, dalam $\Sigma_m$; yaitu, himpunan kelas ekuivalensi dalam
$\Omega$ di bawah relasi $\omega\sim\omega'$ jika
$Z_{kn}(\omega)=Z_{kn}(\omega')$ untuk semua $k$, $n\le m$. Untuk sebarang
variabel acak terintegralkan $X$ pada $\Omega$, definisikan $\Expn_m(X)$
dengan menetapkan

$$\eqalign{\Expn_m(X)(\omega)
&=\Bover1{\mu A}\int_AX\text{ jika }\omega\in A\in\Cal A_m
  \text{ dan }\mu A>0,\cr
&=0\text{ jika }\omega\in A\in\Cal A_m\text{ dan }\mu A=0.\cr}$$

\noindent Maka $\Expn_m(X)$ adalah ekspektasi bersyarat dari $X$ pada
$\Sigma_m$.
%It will be useful later to note that $\Expn_l(\Expn_m(X))=\Expn_l(X)$
%whenever $l\le m$.

Sekarang

$$\eqalignno{\lim_{n\to\Cal F}\int|\Expn_m(F_k(X_n)-Y_k)|
&=\lim_{n\to\Cal F}\sum_{A\in\Cal A_m}\int_A|\Expn_m(F_k(X_n)-Y_k)|\cr
&=\lim_{n\to\Cal F}\sum_{A\in\Cal A_m}|\int_A\Expn_m(F_k(X_n)-Y_k)|\cr
\displaycause{karena $\Expn_m(F_k(X_n)-Y_k)$ konstan pada setiap
$A\in\Cal A_m$}
&=\lim_{n\to\Cal F}\sum_{A\in\Cal A_m}|\int_AF_k(X_n)-Y_k|\cr
&=\sum_{A\in\Cal A_m}\lim_{n\to\Cal F}|\int_AF_k(X_n)-Y_k|
=0\cr}$$

\noindent menurut pemilihan $Y_k$. Jadi jika kita menetapkan

\Centerline{$I_m
=\{n:n\in\Bbb N,\,\int|\Expn_m(F_k(X_n)-Y_k)|\le 2^{-k}$ untuk setiap
$k\le m\}$,}

\noindent maka $I_m\in\Cal F$ untuk setiap $m$.

\medskip

{\bf (d)} Andaikan $\sequencen{r(n)}$ adalah sebarang barisan naik ketat
dalam $\Bbb N$ sedemikian sehingga $r(0)>0$, dan
$r(n)\in J_n\cap K_n$ untuk setiap $n$, serta
$r(n)\in I_{r(n-1)}$ untuk $n\ge 1$. Maka
$\Bover1{n+1}\sum_{i=0}^nX_{r(i)}\to Y$ hampir di mana-mana ketika
$n\to\infty$. \Prf\ Nyatakan $X_{r(n)}$ sebagai

\Centerline{$(X_{r(n)}-F_n(X_{r(n)}))
+(F_n(X_{r(n)})-Y_n-Z_{n,r(n)})+Y_n+Z_{n,r(n)}$}

\noindent untuk setiap $n$. Kita bahas komponen-komponen ini satu per satu:

\medskip

\quad{\bf (i)}

$$\eqalignno{\sum_{n=0}^{\infty}\Pr(X_{r(n)}\ne F_n(X_{r(n)}))
&\le\sum_{n=0}^{\infty}p_n+2^{-n}\cr
\displaycause{karena $r(n)\in K_n$ untuk setiap $n$}
&<\infty\cr}$$

\noindent menurut (a). Tetapi ini berarti bahwa
$X_{r(n)}-F_n(X_{r(n)})\to 0$ hampir di mana-mana, sebab barisan tersebut
akhirnya bernilai nol di hampir setiap titik, dan
$\Bover1{n+1}\sum_{i=0}^n(X_{r(i)}-F_i(X_{r(i)}))\to 0$ hampir di mana-mana,
sekali lagi menurut 273Ca.

\medskip

\quad{\bf (ii)} Menurut pemilihan $Z_{n,r(n)}$,

\Centerline{$\sum_{n=0}^{\infty}\int|F_n(X_{r(n)})-Y_n-Z_{n,r(n)}|
\le\sum_{n=0}^{\infty}2^{-n}$}

\noindent berhingga, sehingga
$F_n(X_{r(n)})-Y_n-Z_{n,r(n)}\to 0$ hampir di mana-mana dan
$\Bover1{n+1}\sum_{i=0}^n(F_i(X_{r(i)})-Y_i-Z_{i,r(i)})\to 0$ hampir di
mana-mana.

\medskip

\quad{\bf (iii)} Menurut 276G, $Y_n\to Y$ hampir di mana-mana dan
$\Bover1{n+1}\sum_{i=0}^nY_i\to Y$ hampir di mana-mana.

\medskip

\quad{\bf (iv)} Kita mengetahui bahwa, untuk setiap $n\ge 1$,
$r(n)\in I_{r(n-1)}$. Jadi (karena $r(n-1)\ge n$),
$\int|\Expn_{r(n-1)}(F_n(X_{r(n)})-Y_n)|\le 2^{-n}$. Tetapi juga

\Centerline{$\int\bigl|\Expn_{r(n-1)}
  \bigl(F_n(X_{r(n)})-Y_n-Z_{n,r(n)}\bigr)\bigr|
\le\int|F_n(X_{r(n)})-Y_n-Z_{n,r(n)}|\le 2^{-n}$}

\noindent menurut 244M dan pemilihan $Z_{n,r(n)}$,

$$\eqalign{\int|\Expn_{r(n-1)}Z_{n,r(n)}|
&=\int\bigl|\Expn_{r(n-1)}(F_n(X_{r(n)})-Y_n)
  -\Expn_{r(n-1)}(F_n(X_{r(n)})-Y_n-Z_{n,r(n)})\bigr|\cr
&\le 2^{-n+1}\cr}$$

\noindent untuk setiap $n$. Karena itu
$\Expn_{r(n-1)}Z_{n,r(n)}\to 0$ hampir di mana-mana.

Di sisi lain,

$$\eqalignno{\sum_{n=0}^{\infty}\Bover1{(n+1)^2}\int Z_{n,r(n)}^2
&\le\sum_{n=0}^{\infty}\Bover1{(n+1)^2}
  \int(F_n(X_{r(n)})-Y_n)^2\cr
&\le\sum_{n=0}^{\infty}\Bover1{(n+1)^2}
  (1+\sum_{j=0}^n(j+1)^2p'_j)\cr
\displaycause{karena $r(n)\in J_n$}
&\le\sum_{n=0}^{\infty}\Bover1{(n+1)^2}
  +\sum_{j=0}^{\infty}(j+1)^2p'_j
    \sum_{n=j}^{\infty}\Bover1{(n+1)^2}\cr
&\le\Bover{\pi^2}6
  +2\sum_{j=0}^{\infty}(j+1)p'_j\cr}$$

\noindent berhingga. (Di sini saya menggunakan taksiran

\Centerline{$\sum_{n=j}^{\infty}\Bover1{(n+1)^2}
\le\sum_{n=j}^{\infty}\Bover2{n+1}-\Bover2{n+2}=\Bover2{j+1}$.)}

\noindent Menurut 276F, yang diterapkan pada
$\sequencen{\Sigma_{r(n)}}$ dan $\sequencen{Z_{n,r(n)}}$,
$\Bover1{n+1}\sum_{i=0}^nZ_{i,r(i)}\to 0$ hampir di mana-mana.

\medskip

\quad{\bf (v)} Dengan menjumlahkan keempat komponen ini, kita melihat
bahwa $\Bover1{n+1}\sum_{i=0}^nX_{r(i)}\to Y$, sebagaimana diklaim.\ \Qed

\medskip

{\bf (e)} Sekarang pilih sebarang barisan naik ketat
$\sequencen{s(n)}$ dalam $\Bbb N$ sedemikian sehingga $s(0)>0$,
$s(n)\in\bigcap_{m\le n}J_m\cap K_n$ untuk setiap $n$, dan
$s(n)\in I_{s(n-1)}$ untuk $n\ge 1$; barisan semacam itu ada karena
$J_m\cap K_m\cap I_{s(n-1)}$ termasuk dalam $\Cal F$, sehingga tak
berhingga, untuk setiap $n\ge 1$ dan $m\in\Bbb N$. Tetapkan
$X'_n=X_{s(n)}$ untuk setiap $n$. Jika $\sequencen{X''_n}$ adalah
subbarisan dari $\sequencen{X'_n}$, maka barisan itu berbentuk
$\sequencen{X_{s(r(n))}}$ untuk suatu barisan naik ketat
$\sequencen{r(n)}$. Dalam hal ini

\Centerline{$s(r(0))\ge s(0)>0$,}

\Centerline{$s(r(n))\in J_{r(n)}\cap K_{r(n)}\subseteq J_n\cap K_n$
untuk setiap $n$,}

\Centerline{$s(r(n))\in I_{s(r(n)-1)}\subseteq I_{s(r(n-1))}$ untuk setiap
$n\ge 1$.}

\noindent Jadi (d) memberi tahu kita bahwa
$\Bover1{n+1}\sum_{i=0}^nX''_i\to Y$ hampir di mana-mana.

\medskip

{\bf (f)} Dengan demikian teorema telah terbukti untuk kasus
$(\Omega,\Sigma,\mu)$ ruang probabilitas. Sekarang andaikan $\mu$
$\sigma$-hingga dan $\mu\Omega>0$. Dalam hal ini terdapat fungsi terukur
positif ketat $f:\Omega\to\Bbb R$ sedemikian sehingga $\int fd\mu=1$
(215B(ix)). Misalkan $\nu$ ukuran integral tak tentu yang bersesuaian
(234J), sehingga $\nu$ adalah ukuran probabilitas pada $\Omega$, dan
$\sequencen{\Bover1f\times X_n}$ suatu barisan fungsi terintegralkan
terhadap $\nu$ sedemikian sehingga
$\sup_{n\in\Bbb N}\int|\Bover1f\times X_n|d\nu$ berhingga (235K). Dari
(a)--(e) kita melihat bahwa harus terdapat suatu fungsi yang terintegralkan
terhadap $\nu$, yaitu $Y$, dan suatu subbarisan $\sequencen{X'_n}$ dari
$\sequencen{X_n}$ sedemikian sehingga
$\Bover1{n+1}\sum_{i=0}^n\Bover1f\times X''_i\to Y\,\,\nu$-hampir di
mana-mana untuk setiap subbarisan $\sequencen{X''_n}$ dari
$\sequencen{X'_n}$. Tetapi $\mu$ dan $\nu$ mempunyai himpunan terabaikan
yang sama (234Lc), sehingga
$\Bover1{n+1}\sum_{i=0}^nX''_i\to f\times Y\,\,\mu$-hampir di mana-mana
untuk setiap subbarisan $\sequencen{X''_n}$ dari $\sequencen{X'_n}$.

\medskip

{\bf (g)} Karena hasilnya trivial jika $\mu\Omega=0$, teorema berlaku
kapan pun $\mu$ adalah $\sigma$-hingga. Untuk kasus umum, tetapkan

\Centerline{$\tilde\Omega
=\bigcup_{n\in\Bbb N}\{\omega:X_n(\omega)\ne 0\}
=\bigcup_{m,n\in\Bbb N}\{\omega:|X_n(\omega)|\ge 2^{-m}\}$,}

\noindent sehingga ukuran subruang $\mu_{\tilde\Omega}$ adalah
$\sigma$-hingga. Maka terdapat suatu fungsi yang terintegralkan terhadap
$\mu_{\tilde\Omega}$, yaitu $\tilde Y$, dan suatu subbarisan
$\sequencen{X'_n}$ dari $\sequencen{X_n}$ sedemikian sehingga
$\Bover1{n+1}\sum_{i=0}^nX''_i\restrp\tilde\Omega\to\tilde Y\,\,
\mu_{\tilde\Omega}$-hampir di mana-mana untuk setiap subbarisan
$\sequencen{X''_n}$ dari $\sequencen{X'_n}$. Dengan menetapkan
$Y(\omega)=\tilde Y(\omega)$ jika $\omega\in\tilde\Omega$, dan $0$ untuk
$\omega\in\Omega\setminus\tilde\Omega$, kita melihat bahwa $Y$
terintegralkan terhadap $\mu$ dan bahwa
$\Bover1{n+1}\sum_{i=0}^nX''_i\to Y\,\,\mu$-hampir di mana-mana kapan
pun $\sequencen{X''_n}$ adalah subbarisan dari $\sequencen{X'_n}$.
Ini menyelesaikan bukti.
}%end of proof of 276H

\exercises{\leader{276X}{Latihan dasar $\pmb{>}$(a)}
%\sqheader 276Xa
Misalkan $\sequencen{X_n}$ adalah martingal yang disesuaikan dengan
suatu barisan $\sequencen{\Sigma_n}$ aljabar-sigma. Tunjukkan bahwa
$\int_EX_n^2\le\int_EX_{n+1}^2$ kapan pun $n\in\Bbb N$ dan
$E\in\Sigma_n$ (dengan memperbolehkan $\infty$ sebagai nilai integral).
\Hint{lihat bukti 276B.}
%276B

\sqheader 276Xb Misalkan $\sequencen{X_n}$ adalah martingal. Tunjukkan
bahwa untuk sebarang $\epsilon>0$,

\Centerline{$\Pr(\sup_{n\in\Bbb N}|X_n|\ge\epsilon)
\le\Bover1{\epsilon^2}\sup_{n\in\Bbb N}\Expn(X_n^2)$.}

\noindent\Hint{gabungkan 276Xa dengan argumen untuk 275D.}
%276Xa, 276B

\spheader 276Xc Kapan 276Xb memberikan hasil yang lebih tajam daripada
275Xb?
%276Xb,  276B

\spheader 276Xd Misalkan $\sequencen{X_n}$ adalah barisan variabel acak
bebas dan berdistribusi identik, dengan ekspektasi nol dan varians hingga
tak nol, dan $\sequencen{t_n}$ suatu barisan dalam $\Bbb R$. Tunjukkan
bahwa (i) jika $\sum_{n=0}^{\infty}t_n^2<\infty$, maka
$\sum_{n=0}^{\infty}t_nX_n$ terdefinisi dalam $\Bbb R$ hampir di
mana-mana (ii) jika $\sum_{n=0}^{\infty}t_n^2=\infty$, maka
$\sum_{n=0}^{\infty}t_nX_n$ tidak terdefinisi hampir di mana-mana.
\Hint{276B, 274Xj.}
%276B

\spheader 276Xe Misalkan $\sequencen{X_n}$ adalah barisan selisih
martingal dan tetapkan $Y_n=\bover1{n+1}(X_0+\ldots+X_n)$ untuk setiap
$n\in\Bbb N$. Tunjukkan bahwa jika $\sequencen{X_n}$ terintegralkan
secara seragam, maka $\lim_{n\to\infty}\|Y_n\|_1=0$. \Hint{gunakan
argumen 273Na, dengan 276C menggantikan 273D, dan tetapkan
$\tilde X_n=X'_n-Z_n$, dengan $Z_n$ ekspektasi bersyarat yang sesuai dari
$X'_n$.}
%276C

\spheader 276Xf Andaikan bahwa $\sequencen{X_n}$ adalah barisan selisih
martingal yang terbatas seragam dan $\sequencen{a_n}\in\ell^2$.
Tunjukkan bahwa $\lim_{n\to\infty}\prod_{i=0}^n(1+a_iX_i)$ terdefinisi
dan berhingga hampir di mana-mana. \Hint{$\sequencen{a_nX_n}$ dapat
dijumlahkan dan dapat dikuadrat-jumlahkan hampir di mana-mana.}
%276C

\sqheader 276Xg\dvAformerly{2{}76Xe}{\bf Hukum kuat bilangan besar:
bentuk kelima} Suatu barisan $\sequencen{X_n}$ variabel acak disebut
{\bf dapat dipertukarkan} jika $(X_{n_0},\ldots,X_{n_k})$ mempunyai
distribusi bersama yang sama dengan $(X_0,\ldots,X_k)$ kapan pun
$n_0,\ldots,n_k$ berbeda-beda. Tunjukkan bahwa jika $\sequencen{X_n}$
adalah barisan variabel acak yang dapat dipertukarkan dan berekspektasi
hingga, maka $\sequencen{\Bover1{n+1}\sum_{i=0}^{n}X_i}$ konvergen
hampir di mana-mana. \Hint{276H.}
%276H

\leader{276Y}{Latihan lanjutan (a)}
%\spheader 276Ya
Misalkan $\sequencen{X_n}$ adalah barisan selisih martingal sedemikian
sehingga $\sup_{n\in\Bbb N}
\ifdim\pagewidth>467pt\penalty-100\fi
\|X_n\|_p$ berhingga, dengan $p\in\ooint{1,\infty}$. Tunjukkan bahwa
$\lim_{n\to\infty}\|\bover1{n+1}\sum_{i=0}^nX_i\|_p=0$.
\Hint{273Nb.}
%276A

\spheader 276Yb Misalkan $\sequencen{X_n}$ adalah barisan selisih
martingal yang terintegralkan secara seragam dan $Y$ suatu variabel acak
terbatas. Tunjukkan bahwa
$\lim_{n\to\infty}\Expn(X_n\times Y)=0$. (Bandingkan 272Ye.)
%276A

\spheader 276Yc Gunakan 275Yh untuk membuktikan 276Xa.
%276Xa, 276B

\spheader 276Yd Misalkan $\sequencen{X_n}$ adalah barisan variabel acak
sedemikian sehingga, untuk suatu $\delta>0$,
$\sup_{n\in\Bbb N}n^{\delta}\Expn(|X_n|)$ berhingga. Tetapkan
$S_n=\bover1{n+1}(X_0+\ldots+X_n)$ untuk setiap $n$. Tunjukkan bahwa
$\lim_{n\to\infty}S_n=0$ hampir di mana-mana. \Hint{tetapkan
$Z_k=2^{-k}(|X_0|+\ldots+|X_{2^k-1}|)$. Tunjukkan bahwa
$\sum_{k=0}^{\infty}\Expn(Z_k)<\infty$. Tunjukkan bahwa
$|S_n|\le 2Z_{k+1}$ jika $2^k<n\le 2^{k+1}$.}
%for 276Ye

\spheader 276Ye {\bf Hukum kuat bilangan besar: bentuk keenam} Misalkan
$\sequencen{X_n}$ adalah barisan selisih martingal sedemikian sehingga,
untuk suatu $\delta>0$,
$\sup_{n\in\Bbb N}\Expn(|X_n|^{1+\delta})$ berhingga. Tetapkan
$S_n=\bover1{n+1}(X_0+\ldots+X_n)$ untuk setiap $n$. Tunjukkan bahwa
$\lim_{n\to\infty}S_n=0$ hampir di mana-mana. \Hint{ambil suatu barisan
tak menurun $\sequencen{\Sigma_n}$ yang menjadi tempat penyesuaian
$\sequencen{X_n}$. Tetapkan $Y_n=X_n$ ketika $|X_n|\le n$, dan $0$ jika
tidak. Misalkan $U_n$ suatu ekspektasi bersyarat dari $Y_n$ pada
$\Sigma_{n-1}$ dan tetapkan $Z_n=Y_n-U_n$. Gunakan gagasan dari 273H,
276C, dan 276Yd di atas untuk menunjukkan bahwa
$\bover{1}{n+1}\sum_{i=0}^nV_i\to 0$ hampir di mana-mana untuk
$V_i=Z_i$, $V_i=U_i$, dan $V_i=X_i-Y_i$.}
%276Yd, 276C mt27bits

\spheader 276Yf Tunjukkan bahwa terdapat martingal $\sequencen{X_n}$
yang konvergen dalam ukuran tetapi tidak konvergen hampir di mana-mana.
(Bandingkan 273Ba.) \Hint{atur agar
$\{\omega:X_{n+1}(\omega)\ne 0\}=E_n
\subseteq\{\omega:|X_{n+1}(\omega)-X_n(\omega)|\ge 1\}$, dengan
$\sequencen{E_n}$ suatu barisan bebas himpunan dan
$\mu E_n=\Bover1{n+1}$ untuk setiap $n$.}
%with 276Yg

\spheader 276Yg Berikan contoh barisan selisih martingal
$\sequencen{X_n}$ yang berdistribusi identik sedemikian sehingga
$\sequencen{\bover1{n+1}(X_0+\ldots+X_n)}$ tidak konvergen ke $0$ hampir
di mana-mana. \Hint{mulailah dengan menyusun suatu barisan terbatas
seragam $\sequencen{U_n}$ sedemikian sehingga
$\lim_{n\to\infty}\Expn(|U_n|)=0$, tetapi
$\sequencen{\bover1{n+1}(U_0+\ldots+U_n)}$ tidak konvergen ke $0$ hampir
di mana-mana. Sekarang ulangi konstruksi Anda dalam konteks sedemikian
rupa sehingga $U_n$ dapat diturunkan dari suatu barisan selisih martingal
yang berdistribusi identik melalui rumus-rumus dalam 276Ye.}
%276Ye, 276C

\spheader 276Yh Susunlah bukti teorema Koml\'os yang tidak melibatkan
ultrafilter, atau penggunaan lain apa pun dari aksioma pilihan penuh,
tetapi seluruhnya berjalan dengan memilih sub-subbarisan yang sesuai.
Ingatlah untuk memeriksa bahwa Anda dapat membuktikan setiap fakta yang
Anda gunakan mengenai barisan yang konvergen lemah dalam $L^1$ dengan
aturan yang sama.
}%end of exercises

\endnotes{
\Notesheader{276}
Saya menyertakan dua bentuk lagi dari hukum kuat bilangan besar (276C,
276Ye), bukan karena saya memikirkan suatu penerapan, melainkan karena
menurut saya, jika Anda mengetahui hukum kuat untuk barisan bebas yang
terbatas dalam $\|\,\|_{1+\delta}$ dan mengetahui apa itu barisan selisih
martingal, maka ada sesuatu yang kurang jika Anda tidak mengetahui hukum
kuat untuk barisan selisih martingal yang terbatas dalam
$\|\,\|_{1+\delta}$. Lalu, tentu saja, saya harus menambahkan 276Yf dan
276Yg (yang tampaknya sulit), agar Anda tidak tergoda untuk menganggap
bahwa hukum kuat itu `sebenarnya' mengenai barisan selisih martingal,
bukan mengenai barisan bebas. (Bandingkan 272Yd dan 275Xl.)

Teorema Koml\'os agak berada di luar cakupan jilid ini; pembuktiannya
cukup berat dan tentunya jauh kurang penting, bagi kebanyakan ahli
probabilitas, daripada banyak hasil yang saya hilangkan. Teorema ini
memang memberikan bukti singkat untuk 276Xg. Namun teorema tersebut
relevan bagi pertanyaan yang muncul dalam beberapa topik yang dibahas
dalam Jilid 3 dan 4, dan buktinya secara alami cocok ditempatkan dalam
bagian ini.
}%end of notes

\discrpage
